\documentclass[11pt]{article}

\usepackage[margin=1in]{geometry}
\usepackage{amsmath,amssymb,amsthm,mathtools}
\usepackage{hyperref}
\usepackage{enumitem}

\newtheorem{theorem}{Theorem}
\newtheorem{corollary}{Corollary}
\newtheorem{proposition}{Proposition}
\newtheorem{lemma}{Lemma}
\newtheorem{assumption}{Assumption}
\newtheorem{remark}{Remark}

\newcommand{\D}{D}
\newcommand{\R}{\mathbb R}
\newcommand{\E}{\mathbb E}
\newcommand{\Law}{\operatorname{Law}}
\newcommand{\esssup}{\operatorname*{ess\,sup}}
\newcommand{\KL}{D_{\mathrm{KL}}}
\newcommand{\dd}{\,\mathrm d}

\title{A Formal Bound on the Suboptimality of zCDP Accounting under Heterogeneous Composition}
\author{}
\date{}

\begin{document}

\maketitle

\section{Problem Setup}

Consider a heterogeneous composition of mechanisms
\[
M=(M_1,\dots,M_n),
\]
where each mechanism uses the same noise family \(\nu\), possibly with different
scales or sensitivities. Suppose that each \(M_i\) satisfies \(\rho_i\)-zCDP, so
that the usual zCDP composition theorem gives the total privacy parameter
\[
\rho=\sum_{i=1}^n \rho_i.
\]
The goal is to quantify the suboptimality of this scalar accounting rule relative
to the exact Renyi privacy curve.

Throughout this note, we use the convention that a mechanism is \(\rho\)-zCDP if,
for every order \(\alpha>1\),
\[
D_\alpha(P\|Q)\le \rho\alpha,
\]
where \(P,Q\) are the output laws of the mechanism on neighboring datasets.
Under this normalization, define the zCDP accounting defect at order \(\alpha\) by
\[
\Gamma_{n,\nu}(\alpha)
:=
\rho\alpha-D_\alpha(P\|Q).
\]
Equivalently, because Renyi divergence tensorizes over product distributions, one
may also write
\[
\Gamma_{n,\nu}(\alpha)
:=
\alpha\sum_{i=1}^n\rho_i-
\sum_{i=1}^n D_\alpha(P_i\|Q_i).
\]
The purpose of the results below is to give lower bounds on
\(\Gamma_{n,\nu}(\alpha)\) as a function of the number of compositions \(n\), the
allocation \((\rho_1,\dots,\rho_n)\), and the noise law \(\nu\).

\section{Assumptions}

We work under the following standing assumptions.

\begin{assumption}[Independent additive composition]
\label{ass:additive-composition}
For each coordinate \(i\), the two output laws of \(M_i\) on a neighboring pair of
datasets are
\[
P_i=\Law(Z_i+\theta_i),
\qquad
Q_i=\Law(Z_i),
\]
where \(Z_i\sim \nu_{s_i}\), and the random variables \(Z_1,\dots,Z_n\) are
independent. The distributions \(\nu_{s_i}\) belong to a common noise family
\(\nu\), possibly with different scales \(s_i\). Thus the joint output laws are
\[
P=\bigotimes_{i=1}^n P_i,
\qquad
Q=\bigotimes_{i=1}^n Q_i.
\]
\end{assumption}

\begin{assumption}[Absolute continuity and finite Renyi divergences]
\label{ass:regularity}
For every coordinate \(i\), assume
\[
P_i\ll Q_i,
\qquad
Q_i\ll P_i,
\]
and
\[
D_\alpha(P_i\|Q_i)<\infty
\qquad
\text{for all } \alpha>1.
\]
\end{assumption}

\begin{remark}
The absolute-continuity and finiteness assumptions hold for many smooth additive
noise families with everywhere positive densities. Centering and symmetry of
\(\nu\) are not needed for the tensorization identity, but they are useful when
one wants local expansions of the per-coordinate Renyi curve in terms of the
shift size. In particular, symmetry typically removes odd-order terms in the
small-shift expansion.
\end{remark}

\section{Exact Tensorization of the Renyi Curve}

\begin{lemma}[Tensorization of Renyi divergence]
\label{lem:tensorization}
Under Assumptions~\ref{ass:additive-composition} and~\ref{ass:regularity}, for
every \(\alpha>1\),
\[
D_\alpha(P\|Q)
=
\sum_{i=1}^n D_\alpha(P_i\|Q_i).
\]
\end{lemma}

\begin{proof}
By definition,
\[
D_\alpha(P\|Q)
=
\frac{1}{\alpha-1}
\log \int
\left(\frac{\dd P}{\dd Q}\right)^\alpha \dd Q.
\]
Since \(P=\bigotimes_i P_i\) and \(Q=\bigotimes_i Q_i\), we have
\[
\frac{\dd P}{\dd Q}(x_1,\dots,x_n)
=
\prod_{i=1}^n \frac{\dd P_i}{\dd Q_i}(x_i).
\]
Therefore,
\[
\begin{aligned}
\int
\left(\frac{\dd P}{\dd Q}\right)^\alpha \dd Q
&=
\int
\prod_{i=1}^n
\left(\frac{\dd P_i}{\dd Q_i}(x_i)\right)^\alpha
\prod_{i=1}^n \dd Q_i(x_i) \\
&=
\prod_{i=1}^n
\int
\left(\frac{\dd P_i}{\dd Q_i}\right)^\alpha \dd Q_i.
\end{aligned}
\]
Taking logarithms and dividing by \(\alpha-1\) gives
\[
D_\alpha(P\|Q)
=
\sum_{i=1}^n D_\alpha(P_i\|Q_i),
\]
as claimed.
\end{proof}

\section{The Exact zCDP Accounting Defect}

For each coordinate define its exact zCDP radius by
\[
\bar\rho_i
:=
\sup_{\beta>1}\frac{D_\beta(P_i\|Q_i)}{\beta}.
\]
If a mechanism is certified using a possibly non-minimal zCDP parameter
\(\rho_i\), then necessarily
\[
\rho_i\ge \bar\rho_i.
\]

\begin{theorem}[Exact zCDP accounting defect under heterogeneous composition]
\label{thm:zcdp-defect-composition}
Under Assumptions~\ref{ass:additive-composition} and~\ref{ass:regularity}, let
\(P=\bigotimes_i P_i\) and \(Q=\bigotimes_i Q_i\). Suppose each coordinate is
certified as \(\rho_i\)-zCDP, where \(\rho_i\ge \bar\rho_i\). Then for every
\(\alpha>1\),
\[
\Gamma_{n,\nu}(\alpha)
:=
\alpha\sum_{i=1}^n\rho_i-D_\alpha(P\|Q)
\]
satisfies the exact identity
\[
\Gamma_{n,\nu}(\alpha)
=
\alpha\sum_{i=1}^n(\rho_i-\bar\rho_i)
+
\sum_{i=1}^n
\Bigl\{\alpha\bar\rho_i-D_\alpha(P_i\|Q_i)\Bigr\}.
\]
In particular,
\[
\Gamma_{n,\nu}(\alpha)\ge 0.
\]
If each coordinate is calibrated to its exact zCDP radius, namely
\(\rho_i=\bar\rho_i\), then
\[
\Gamma_{n,\nu}(\alpha)
=
\sum_{i=1}^n
\Bigl\{\alpha\rho_i-D_\alpha(P_i\|Q_i)\Bigr\}.
\]
Moreover, \(\Gamma_{n,\nu}(\alpha)=0\) for every allocation
\((\rho_1,\dots,\rho_n)\) and every \(\alpha>1\) if every coordinate has a linear
Gaussian-type Renyi curve
\[
D_\alpha(P_i\|Q_i)=\alpha\rho_i
\qquad
\text{for all } \alpha>1.
\]
\end{theorem}

\begin{proof}
By Lemma~\ref{lem:tensorization},
\[
D_\alpha(P\|Q)
=
\sum_{i=1}^n D_\alpha(P_i\|Q_i).
\]
Therefore,
\[
\begin{aligned}
\Gamma_{n,\nu}(\alpha)
&=
\alpha\sum_{i=1}^n \rho_i
-
\sum_{i=1}^n D_\alpha(P_i\|Q_i) \\
&=
\alpha\sum_{i=1}^n(\rho_i-\bar\rho_i)
+
\sum_{i=1}^n
\Bigl\{\alpha\bar\rho_i-D_\alpha(P_i\|Q_i)\Bigr\}.
\end{aligned}
\]
Since
\[
\bar\rho_i
=
\sup_{\beta>1}\frac{D_\beta(P_i\|Q_i)}{\beta},
\]
we have, for every fixed \(\alpha>1\),
\[
D_\alpha(P_i\|Q_i)
\le
\alpha\bar\rho_i.
\]
Thus each term
\[
\alpha\bar\rho_i-D_\alpha(P_i\|Q_i)
\]
is nonnegative. Also \(\rho_i-\bar\rho_i\ge 0\). Hence
\[
\Gamma_{n,\nu}(\alpha)\ge 0.
\]
If \(\rho_i=\bar\rho_i\), the first sum vanishes, giving
\[
\Gamma_{n,\nu}(\alpha)
=
\sum_{i=1}^n
\Bigl\{\alpha\rho_i-D_\alpha(P_i\|Q_i)\Bigr\}.
\]
Finally, if
\[
D_\alpha(P_i\|Q_i)=\alpha\rho_i
\qquad
\text{for all } i \text{ and all } \alpha>1,
\]
then each summand is zero, so
\[
\Gamma_{n,\nu}(\alpha)=0
\]
for every \(\alpha>1\). This is the Gaussian-type equality case. Indeed, for an
additive Gaussian mechanism with shift sensitivity \(\Delta_i\) and variance
\(\sigma_i^2\), one has
\[
D_\alpha(P_i\|Q_i)
=
\frac{\alpha\Delta_i^2}{2\sigma_i^2}
=
\alpha\rho_i,
\]
where
\[
\rho_i=\frac{\Delta_i^2}{2\sigma_i^2}.
\]
\end{proof}

\section{Noise-Dependent Per-Coordinate Defect}

For a fixed noise family \(\nu\), define the per-coordinate envelope defect at
Renyi order \(\alpha>1\) by
\[
g_\nu(\alpha;r)
:=
\alpha r-D_\alpha(P_{\nu,r}\|Q_{\nu,r}),
\]
where \((P_{\nu,r},Q_{\nu,r})\) denotes a neighboring pair calibrated so that the
exact zCDP radius is \(r\). With this notation, under exact calibration,
\[
\Gamma_{n,\nu}(\alpha)
=
\sum_{i=1}^n g_\nu(\alpha;\rho_i).
\]
Thus the accounting defect is additive across coordinates, but the summand is
not linear in \(\rho_i\). The dependence on the allocation is therefore governed
by the shape of \(r\mapsto g_\nu(\alpha;r)\).

\begin{assumption}[Local non-Gaussian envelope defect]
\label{ass:local-defect}
Fix \(\alpha>1\). Assume that the noise family \(\nu\) has a calibrated
per-coordinate defect satisfying the following local expansion: there exist
constants \(r_0>0\), \(c_\nu(\alpha)>0\), and \(C_\nu(\alpha)<\infty\) such that,
for every \(0\le r\le r_0\),
\[
g_\nu(\alpha;r)
\ge
c_\nu(\alpha)r^2-C_\nu(\alpha)r^3.
\]
Equivalently, for sufficiently small \(r\),
\[
g_\nu(\alpha;r)
\gtrsim_\nu r^2.
\]
The constant \(c_\nu(\alpha)\) is the leading non-Gaussian curvature defect of
the Renyi curve relative to its zCDP envelope. For Gaussian noise,
\[
g_\nu(\alpha;r)\equiv 0,
\]
so \(c_\nu(\alpha)=0\).
\end{assumption}

\begin{remark}[Interpretation through privacy-loss cumulants]
Let
\[
L_t(x)
=
\log\frac{\dd P_t}{\dd Q_t}(x)
\]
be the privacy-loss random variable, and let
\[
K_t(\lambda)
=
\log \E_{Q_t}\exp\{\lambda L_t\}
\]
be its cumulant generating function. Then
\[
D_\alpha(P_t\|Q_t)
=
\frac{K_t(\alpha)}{\alpha-1}.
\]
The exact zCDP radius is
\[
r(t)
=
\sup_{\beta>1}
\frac{K_t(\beta)}{\beta(\beta-1)}.
\]
Therefore
\[
g_\nu(\alpha;r(t))
=
\alpha r(t)-\frac{K_t(\alpha)}{\alpha-1}.
\]
Assumption~\ref{ass:local-defect} states that, after calibration by the exact
zCDP radius, the cumulant generating function lies strictly below its best
quadratic zCDP envelope by an amount of order at least \(r^2\). Gaussian privacy
loss is the exceptional case in which the cumulant generating function is exactly
quadratic and the defect vanishes.
\end{remark}

\section{Allocation-Sensitive Lower Bounds}

\begin{corollary}[Allocation-sensitive lower bound]
\label{cor:allocation-lower-bound}
Suppose the assumptions of Theorem~\ref{thm:zcdp-defect-composition} hold and
each coordinate is calibrated to its exact zCDP radius:
\[
\rho_i=\bar\rho_i.
\]
Assume further that Assumption~\ref{ass:local-defect} holds at the fixed order
\(\alpha>1\), and that
\[
0\le \rho_i\le r_*
:=
\min\left\{r_0,\frac{c_\nu(\alpha)}{2C_\nu(\alpha)}\right\}
\]
for every \(i\). Then
\[
\Gamma_{n,\nu}(\alpha)
\ge
\frac{c_\nu(\alpha)}{2}
\sum_{i=1}^n \rho_i^2.
\]
Equivalently, writing
\[
\rho=\sum_{i=1}^n \rho_i,
\qquad
w_i=\frac{\rho_i}{\rho},
\qquad
H(w)=\sum_{i=1}^n w_i^2,
\]
one has
\[
\Gamma_{n,\nu}(\alpha)
\ge
\frac{c_\nu(\alpha)}{2}\rho^2H(w).
\]
Since
\[
H(w)\ge \frac1n,
\]
the allocation minimizing the guaranteed defect for fixed total \(\rho\) is the
uniform allocation
\[
\rho_i=\frac{\rho}{n}.
\]
In that case,
\[
\Gamma_{n,\nu}(\alpha)
\ge
\frac{c_\nu(\alpha)}{2}\frac{\rho^2}{n}.
\]
More generally, defining the effective number of compositions by
\[
n_{\mathrm{eff}}
:=
\frac{\rho^2}{\sum_{i=1}^n\rho_i^2}
=
\frac1{H(w)},
\]
one obtains
\[
\Gamma_{n,\nu}(\alpha)
\ge
\frac{c_\nu(\alpha)}{2}
\frac{\rho^2}{n_{\mathrm{eff}}}.
\]
\end{corollary}

\begin{proof}
By Theorem~\ref{thm:zcdp-defect-composition}, under exact calibration,
\[
\Gamma_{n,\nu}(\alpha)
=
\sum_{i=1}^n g_\nu(\alpha;\rho_i).
\]
By Assumption~\ref{ass:local-defect},
\[
g_\nu(\alpha;\rho_i)
\ge
c_\nu(\alpha)\rho_i^2-C_\nu(\alpha)\rho_i^3.
\]
Since
\[
\rho_i\le \frac{c_\nu(\alpha)}{2C_\nu(\alpha)},
\]
we have
\[
C_\nu(\alpha)\rho_i^3
\le
\frac{c_\nu(\alpha)}{2}\rho_i^2.
\]
Therefore,
\[
g_\nu(\alpha;\rho_i)
\ge
\frac{c_\nu(\alpha)}{2}\rho_i^2.
\]
Summing over \(i\) yields
\[
\Gamma_{n,\nu}(\alpha)
\ge
\frac{c_\nu(\alpha)}{2}
\sum_{i=1}^n \rho_i^2.
\]
Now write \(\rho_i=\rho w_i\), where \(\sum_i w_i=1\). Then
\[
\sum_{i=1}^n \rho_i^2
=
\rho^2\sum_{i=1}^n w_i^2
=
\rho^2H(w).
\]
By Cauchy--Schwarz,
\[
\left(\sum_{i=1}^n w_i\right)^2
\le
n\sum_{i=1}^n w_i^2.
\]
Since \(\sum_i w_i=1\), this gives
\[
H(w)=\sum_i w_i^2\ge \frac1n.
\]
Equality holds if and only if
\[
w_1=\cdots=w_n=\frac1n,
\]
equivalently
\[
\rho_1=\cdots=\rho_n=\frac{\rho}{n}.
\]
Hence, for the uniform allocation,
\[
\Gamma_{n,\nu}(\alpha)
\ge
\frac{c_\nu(\alpha)}{2}\frac{\rho^2}{n}.
\]
Finally, the effective-composition form follows from
\[
\sum_{i=1}^n \rho_i^2
=
\frac{\rho^2}{n_{\mathrm{eff}}}.
\]
\end{proof}

\begin{proposition}[Convex envelope lower bound]
\label{prop:convex-envelope}
Assume that, for the noise family \(\nu\) and fixed order \(\alpha>1\), the
function
\[
r\mapsto g_\nu(\alpha;r)
\]
is convex on an interval containing \(\rho_1,\dots,\rho_n\). Then, for fixed total
\(\rho=\sum_i\rho_i\),
\[
\Gamma_{n,\nu}(\alpha)
=
\sum_{i=1}^n g_\nu(\alpha;\rho_i)
\ge
n g_\nu\left(\alpha;\frac{\rho}{n}\right).
\]
If, in addition,
\[
g_\nu(\alpha;r)
\ge
 a_\nu(\alpha)r^p
\]
for some \(p\ge 1\), then
\[
\Gamma_{n,\nu}(\alpha)
\ge
 a_\nu(\alpha)
\sum_{i=1}^n \rho_i^p.
\]
In particular, for \(p=2\),
\[
\Gamma_{n,\nu}(\alpha)
\ge
 a_\nu(\alpha)\frac{\rho^2}{n}.
\]
\end{proposition}

\begin{proof}
The identity
\[
\Gamma_{n,\nu}(\alpha)=\sum_{i=1}^n g_\nu(\alpha;\rho_i)
\]
follows from Theorem~\ref{thm:zcdp-defect-composition} under exact calibration.
If \(g_\nu(\alpha;\cdot)\) is convex, Jensen's inequality gives
\[
\frac1n\sum_{i=1}^n g_\nu(\alpha;\rho_i)
\ge
 g_\nu\left(\alpha;\frac1n\sum_{i=1}^n \rho_i\right)
=
 g_\nu\left(\alpha;\frac{\rho}{n}\right).
\]
Multiplying by \(n\) gives
\[
\Gamma_{n,\nu}(\alpha)
\ge
n g_\nu\left(\alpha;\frac{\rho}{n}\right).
\]
If
\[
g_\nu(\alpha;r)
\ge
 a_\nu(\alpha)r^p,
\]
then summing over \(i\) gives
\[
\Gamma_{n,\nu}(\alpha)
\ge
 a_\nu(\alpha)\sum_{i=1}^n \rho_i^p.
\]
For \(p=2\), Cauchy--Schwarz gives
\[
\sum_{i=1}^n \rho_i^2
\ge
\frac1n\left(\sum_{i=1}^n \rho_i\right)^2
=
\frac{\rho^2}{n}.
\]
Thus
\[
\Gamma_{n,\nu}(\alpha)
\ge
 a_\nu(\alpha)\frac{\rho^2}{n}.
\]
\end{proof}

\section{Loose Per-Coordinate Calibration}

The previous bounds separate the intrinsic loss caused by replacing the full
Renyi curve by a scalar zCDP radius from the extrinsic loss caused by certifying a
coordinate with a loose zCDP parameter.

If \(\rho_i\ge \bar\rho_i\), then Theorem~\ref{thm:zcdp-defect-composition} gives
\[
\Gamma_{n,\nu}(\alpha)
=
\alpha\sum_{i=1}^n(\rho_i-\bar\rho_i)
+
\sum_{i=1}^n
\Bigl\{\alpha\bar\rho_i-D_\alpha(P_i\|Q_i)\Bigr\}.
\]
Under the local defect assumption applied to \(\bar\rho_i\), this yields
\[
\Gamma_{n,\nu}(\alpha)
\ge
\alpha\sum_{i=1}^n(\rho_i-\bar\rho_i)
+
\frac{c_\nu(\alpha)}{2}
\sum_{i=1}^n \bar\rho_i^2,
\]
provided \(\bar\rho_i\le r_*\) for all \(i\). The first term is calibration loss:
it is present only because the certified parameter \(\rho_i\) is larger than the
exact zCDP radius \(\bar\rho_i\). The second term is intrinsic to the noise law
and measures the departure of the exact Renyi curve from the Gaussian linear
curve.

\section{Conclusion}

For an independently composed additive mechanism, the exact Renyi curve satisfies
\[
D_\alpha(P\|Q)
=
\sum_{i=1}^nD_\alpha(P_i\|Q_i).
\]
Therefore, under exact per-coordinate zCDP calibration,
\[
\boxed{
\Gamma_{n,\nu}(\alpha)
=
\sum_{i=1}^n
\Bigl\{\alpha\rho_i-D_\alpha(P_i\|Q_i)\Bigr\}.
}
\]
For Gaussian noise,
\[
D_\alpha(P_i\|Q_i)=\alpha\rho_i,
\]
so
\[
\Gamma_{n,\nu}(\alpha)=0
\]
for every number of compositions \(n\), every allocation
\((\rho_1, \dots,\rho_n)\), and every order \(\alpha>1\). Thus Gaussian zCDP
accounting is exactly tight.

For a non-Gaussian smooth symmetric additive noise family satisfying the local
curvature-defect assumption,
\[
\boxed{
\Gamma_{n,\nu}(\alpha)
\gtrsim_\nu
\sum_{i=1}^n\rho_i^2
=
\rho^2\sum_{i=1}^n\left(\frac{\rho_i}{\rho}\right)^2.
}
\]
Equivalently,
\[
\boxed{
\Gamma_{n,\nu}(\alpha)
\gtrsim_\nu
\frac{\rho^2}{n_{\mathrm{eff}}},
\qquad
n_{\mathrm{eff}}
=
\frac{\rho^2}{\sum_i\rho_i^2}.
}
\]
The uniform allocation \(\rho_i=\rho/n\) minimizes the guaranteed defect and gives
\[
\Gamma_{n,\nu}(\alpha)
\gtrsim_\nu
\frac{\rho^2}{n}.
\]
Highly nonuniform allocations decrease \(n_{\mathrm{eff}}\), increase
\(\sum_i\rho_i^2\), and therefore increase the suboptimality of scalar zCDP
accounting relative to the exact Renyi privacy curve.

\end{document}
